\clearpage
\VerdierPhysicalPageAnchor{0226}%
\begingroup\small\setlength{\parskip}{0pt}\setlength{\emergencystretch}{1.5em}
\setlength{\abovedisplayskip}{4pt plus 2pt minus 2pt}%
\setlength{\belowdisplayskip}{4pt plus 2pt minus 2pt}%

\VerdierSemanticAnchor{ch4d-4-1-2-cont}
これらのモノ射は $Y$ の $\widetilde{\mathbb Z}$ 型フィルトレーションを定める。
これを\emph{$Y$ の第一標準フィルトレーション}という。したがって、
$\widetilde{\mathbb Z}$ から $\operatorname{comp}(\mathcal A)$ への関手が定まり、
これも $Y$ と書く。(4.1.1) で思い出した構成から、$K(\mathcal A)$ に値をとる
スペクトル対象 $(\mathcal Y(p,q),\delta(p,q,r))$ が得られる。ここでは
$(\mathcal Y(p,q),\delta(p,q,r))$ と同型なスペクトル対象を記述する。
まず次のように置く：

\VerdierSemanticAnchor{ch4d-4-1-2-3}
\begin{equation}
\resizebox{0.98\textwidth}{!}{%
$\displaystyle
\begin{aligned}
\text{a) }& -\infty<p<q<+\infty:\\[-2pt]
&Y(p,q)=\cdots\longrightarrow0\longrightarrow0\longrightarrow Y^{-q}
\longrightarrow Y^{-q+1}\longrightarrow\cdots\longrightarrow Y^{-p-2}
\longrightarrow Y^{-p-1}\longrightarrow0\longrightarrow0
\longrightarrow\cdots;\\[2pt]
\text{b) }& -\infty<p<+\infty:\\[-2pt]
&Y(p,+\infty)=\cdots\longrightarrow Y^i\longrightarrow Y^{i+1}
\longrightarrow\cdots\longrightarrow Y^{-p-2}\longrightarrow Y^{-p-1}
\longrightarrow0\longrightarrow0\longrightarrow\cdots;\\[2pt]
\text{c) }&Y(-\infty,q)=Y(q)\quad(q\in\widetilde{\mathbb Z}),\qquad
Y(p,p)=0\quad(p\in\widetilde{\mathbb Z}).
\end{aligned}$%
}
\tag{4.1.2.3}
\end{equation}
これらの複体の微分は複体 $Y$ の微分から誘導される。

次のように書く：

\VerdierSemanticAnchor{ch4d-4-1-2-4}
\begin{equation}
\varphi\bigl((p,q),(p',q')\bigr)\colon
Y(p,q)\longrightarrow Y(p',q'),\qquad p\leq p',\quad q\leq q',
\tag{4.1.2.4}
\end{equation}
これは次の複体の射のホモトピー類である：

\VerdierSemanticAnchor{ch4d-4-1-2-5}
\begin{equation}
\psi\bigl((p,q),(p',q')\bigr)\colon
Y(p,q)\longrightarrow Y(p',q'),
\tag{4.1.2.5}
\end{equation}
その各成分は、始域と終域が非零対象であるとき恒等射である。こうして
$\mathcal{F}\ell_1(\widetilde{\mathbb Z})$ から $K(\mathcal A)$ への関手が
定義されたことは直ちに確かめられる。

$(p,q,r)$ を $\mathcal{F}\ell_2(\widetilde{\mathbb Z})$ の元とする。
次のように書く：

\VerdierSemanticAnchor{ch4d-4-1-2-6}
\begin{equation}
\delta(p,q,r)\colon Y(q,r)\longrightarrow Y(p,q)[1]
\tag{4.1.2.6}
\end{equation}
これは次の複体の射のホモトピー類である：

\VerdierSemanticAnchor{ch4d-4-1-2-7}
\begin{equation}
d(p,q,r)\colon Y(q,r)\longrightarrow Y(p,q)[1],
\tag{4.1.2.7}
\end{equation}
その成分は

\VerdierSemanticAnchor{ch4d-4-1-2-8}
\begin{equation}
\begin{aligned}
d(p,q,r)^i&=0,\qquad i\neq-q-1,\\
d(p,q,r)^{-q-1}&=-d_Y^{-q-1}.
\end{aligned}
\tag{4.1.2.8}
\end{equation}
\endgroup
